


%-------------------------------------------------------------------------

\section{Related Work}
\label{sec:related_work}

\noindent \textbf{Video Generation and Controllable Editing.}
\label{subsec:video_editing}
Diffusion models have significantly advanced video generation~\cite{singer2022make,blattmann2023stable,yang2024cogvideox, ma2026group, ma2024followpose, ma2025followcreation, ma2026fastvmt, ma2025followyourmotion} and editing. To achieve versatile control~\cite{ma2025controllable}, unified frameworks like VACE~\cite{jiang2025vace} propose all-in-one architectures, while EditVerse~\cite{ju2025editverse} and UNIC~\cite{ye2025unic} employ in-context learning to handle diverse tasks. These approaches typically integrate source videos and dense multi-modal conditions into massive joint representations or extended context sequences. Meanwhile, models like InsV2V~\cite{cheng2023consistent} and Lucy Edit~\cite{decart2025lucyedit} leverage synthetic datasets or channel-wise concatenation for high-fidelity modifications~\cite{mokady2023null}. Additionally, reinforcement learning has been explored to further align generated content with human intents~\cite{black2023training}.

Despite their impressive visual quality, these methods inherently rely on offline, non-causal processing. Processing entire video frames and complex conditions jointly quadratically increases the computational overhead. Consequently, these models must observe the entire temporal context before outputting the initial frame. This unacceptable latency intrinsically hinders their deployment in real-time augmented reality scenarios. To break this paradigm, we propose a novel streaming video editing framework explicitly designed for extreme real-time responsiveness.

\noindent \textbf{Streaming Autoregressive Video Models.}
\label{subsec:streaming_models}
Streaming autoregressive generation effectively overcomes the latency bottlenecks of offline models. Foundational works like Diffusion Forcing~\cite{chen2024diffusion} and StreamDiffusion~\cite{kodaira2025streamdiffusion} redefine denoising as block-wise sequential processing. This paradigm enables massive-scale world models (MAGI-1~\cite{teng2025magi}) and infinite-length cinematic generation (SkyReels-V2~\cite{chen2025skyreels}, StreamingT2V~\cite{henschel2025streamingt2v}). To enhance stability and efficiency, Rolling Forcing~\cite{liu2025rolling} suppresses error accumulation, Stable Video Infinity~\cite{li2025stable} mitigates autoregressive drift for infinite-length generation via error-recycling fine-tuning, Self Forcing~\cite{huang2025self} bridges exposure bias via self-generated conditioning, and StreamDiffusionV2~\cite{feng2025streamdiffusionv2} introduces sink-token-guided rolling KV caches for ultra-low-latency live generation. Furthermore, EgoEdit~\cite{li2025egoedit} pioneered streaming models for egocentric video editing by maintaining temporal consistency in continuous first-person visual translations.

However, directly migrating these generation-tailored mechanisms to general video editing introduces a fundamental misalignment. Video generation synthesizes motion from scratch, heavily relying on past generated predictions (e.g., Self Forcing's serial feedback). Conversely, streaming editing is strongly conditioned on continuous source video streams, prioritizing precise spatial restoration over free-form motion synthesis. Serially depending on historical outputs or employing generation-oriented caches (e.g., DeepCache~\cite{ma2024deepcache}, VMem~\cite{li2025vmem}, or StreamDiffusionV2~\cite{feng2025streamdiffusionv2}) degrades high-frequency structural details when rigorously aligning with semantic editing trajectories. Discarding redundant serial feedback, we introduce an AR-based Cache that achieves extreme computational compression while ensuring strict temporal consistency and spatial alignment.

\noindent \textbf{Efficiency and Distillation in Diffusion.}
\label{subsec:distillation}
Accelerating reverse diffusion~\cite{song2023consistencymodels,wang2023videolcm,salimans2022progressive, zheng2025compute, liu2025survey, zheng2026forecast} is crucial for real-time performance. Early one-step distillation efforts focused on the image domain, utilizing techniques like Rectified Flow (InstaFlow~\cite{liu2023instaflow}), target score distillation (TSD-SR~\cite{Dong_2025_CVPR}), and Distribution Matching Distillation~\cite{yin2024improved,yin2025slow}, alongside progressive adversarial distillation (SDXL-Lightning~\cite{lin2024sdxl}). These advancements have naturally extended to the video domain, enabling efficient one-step generation (AAPT~\cite{lin2025diffusion}), restoration (SeedVR2~\cite{wang2025seedvr2}), and super-resolution (DOVE~\cite{chen2025doveefficientonestepdiffusion}). Sparse VideoGen2~\cite{yang2025sparsevideogen2acceleratevideo} further reduces redundancy via attention permutation. Orthogonal to distillation, structural optimizations like Token Merging~\cite{bolya2022token} and FlashAttention~\cite{dao2022flashattention} are widely adopted to alleviate system memory overheads.

Most closely related to our work is FlashVSR~\cite{zhuang2025flashvsr}, which achieves real-time streaming video super-resolution via a three-stage distillation pipeline designed to progressively accelerate latent rendering. Similarly, PersonaLive~\cite{li2025personalive} leverages appearance distillation for live portrait animation. However, these tasks fundamentally involve low-level pixel mapping or structurally constrained regions. Directly transferring these distillation schemes to general video editing—which involves complex semantic reconstruction, cross-scale feature evolution, and training-free reward-guided control—frequently degrades high-fidelity details. In contrast, we propose a tailored three-stage distillation strategy integrated with our AR-based Cache, achieving unprecedented inference acceleration without sacrificing complex editing fidelity.